
\subsubsection{Interpreting Quality-Gated Diversity}

The central question is: why does quality-first (QD) win even when diversity coefficient dominates at large scale ($\alpha _{D}=0.74$ vs $\alpha _{Q}=0.36$ at 10M tokens)?

Our coefficient analysis shows that at 10M tokens, diversity coefficient exceeds quality coefficient. If quality and diversity contributions were independent and additive, diversity-first (DQ) should win when $\alpha _D$ > $\alpha _Q$. The observed QD superiority suggests a quality-gated diversity interaction: diversity sampling appears most effective on quality-filtered subsets. When applied to raw noisy corpora, diversity metrics may capture uninformative variation. Quality filtering removes this noise first, allowing diversity sampling to select among genuinely informative diverse examples.

Consider the signal processing analogy: quality filtering is denoising, diversity sampling is equalization. Equalizing before denoising amplifies noise across frequency bands. Denoising before equalizing preserves signal-to-noise ratio while enhancing pattern coverage.

\subsubsection{Connection to Prior Work}

\textbf{DATAMASK}: Made qualitative observations about quality saturation and diversity persistence. We provide statistical rigor (regression slopes, p-values, effect sizes) and extend with ordering effects not tested by DATAMASK.

\textbf{Data Mixing Laws}: Assume data sources mix independently when combining domain proportions. We show that order-invariance does not hold for sequential filtering steps (Cohen's d=0.76--2.48), extending the framework from domain mixing to sequential filtering contexts.

\textbf{V-Information and FAC}: V-Info showed quality reduction maintains performance at <1M scale. We extend to saturation regime (10M tokens). FAC-Synthesis validated $\rho$=0.90 correlation but did not test scale dependence or ordering. We show evidence that diversity sampling effectiveness may depend on operating post-quality-filtering.

\subsubsection{Limitations}

\textbf{Proxy model size} (GPT-2 Small, 124M parameters): Findings are most applicable to medium-scale model training (100M--1B parameters). Validation with Llama 3 8B is needed for frontier model applicability.

\textbf{Single data source} (C4 realnewslike): Generalization to multi-domain corpora (Pile, RedPajama) or domain-specific data (medical, code) is unknown.

\textbf{Fixed reduction rates} (70\%$\rightarrow 49$\%): Optimal rates likely vary by scale. Adaptive schedules merit investigation.

\textbf{Simulated diversity trajectory} (H-E2): Due to environment constraints, diversity persistence results are modeled from FAC literature rather than direct experiments. Real h-e2 experiments planned for full validation. This does not invalidate QD vs DQ ordering findings (H-M1), which use real training runs.

\textbf{Metric orthogonality} ($\rho$=0.23, p=0.08): While correlation magnitude is low, marginal p-value means we cannot definitively rule out some shared variance between V-Info and FAC. Larger sample validation needed.

\textbf{Coefficient model assumptions}: Compositional model uses simplified linear additive assumptions. More sophisticated models incorporating interaction terms may better capture quality-gating dynamics.

\textbf{Evaluation benchmarks}: MMLU, BEIR, and GSM8K may have contamination overlap with C4 training data. Standard few-shot/zero-shot protocols minimize but do not eliminate this risk.

\textbf{Statistical power}: 3 random seeds per condition provide p<0.05 significance but limited power for detecting small effects. Larger-scale experiments should use 5-10 seeds.

\textbf{Scales tested}: 10K--10M tokens cover practical small-to-medium training regimes but do not reach frontier scale (50M--100M tokens). Reversal may occur beyond tested range.

\subsubsection{Broader Impact}

\textbf{For Practitioners}: For 100M--1B parameter models trained on web corpora at 10K--10M scale, apply quality filtering before diversity sampling. Two-stage curation (QD) requires ~7 GPU-hours for 10M tokens but yields +1.5pp to +5.0pp performance gain.

\textbf{For Research Community}: Compositional ordering effects extend order-invariance assumptions from domain mixing to sequential filtering steps. Trajectory quantification framework (regression slopes, effect sizes) enables principled comparison across studies.

\textbf{Dual-Use Considerations}: Better curation reduces training compute waste but could amplify biases if quality and diversity metrics favor dominant narratives. Practitioners should audit curation metrics for fairness and representativeness.

\subsubsection{Honest Assessment}

We set out to prove scale-dependent reversal: QD dominates at small scale, DQ dominates at large scale. Evidence refutes reversal at 10M tokens but validates underlying mechanisms (quality saturation, modeled diversity persistence, compositional interaction). This pivot from ''when to switch $QD\rightarrow DQ$'' to ''quality-first appears universally beneficial'' is scientifically honest and practically more useful.
